\documentclass[a4paper,13pt]{article}

\usepackage[UTF8,fontset=fandol]{ctex}
\usepackage{geometry}
\usepackage{titlesec}
\usepackage{enumitem}
\usepackage[hidelinks]{hyperref}
\usepackage{fontawesome}
\usepackage{xcolor}
\usepackage{tabularx}

\geometry{left=0.78in,top=0.44in,right=0.78in,bottom=0.5in}

\newcommand{\namefontsize}{\fontsize{22}{26}\selectfont}
\newcommand{\sectiontitlefontsize}{\fontsize{16}{19}\selectfont}
\newcommand{\entrytitlefontsize}{\fontsize{14}{19}\selectfont}
\newcommand{\largeentrytitlefontsize}{\fontsize{15}{19}\selectfont}
\newcommand{\smallentrytitlefontsize}{\fontsize{12}{16}\selectfont}
\newcommand{\entrydetailsfontsize}{\fontsize{11.5}{14}\selectfont}
\newcommand{\contactfontsize}{\fontsize{11}{14}\selectfont}
\newcommand{\linkfontsize}{\fontsize{11.8}{14.6}\selectfont}
\newcommand{\descfontsize}{\fontsize{11}{14.6}\selectfont}
\newcommand{\itemfontsize}{\fontsize{11}{14.6}\selectfont}

\definecolor{applegray}{RGB}{99,99,102}
\definecolor{textgray}{RGB}{85,85,85}
\definecolor{sectiongray}{RGB}{110,110,115}
\definecolor{appledark}{RGB}{0,0,0}
\definecolor{appleblue}{RGB}{0,0,139}

\pagestyle{empty}
\urlstyle{same}
\raggedbottom
\raggedright
\setlength{\tabcolsep}{0in}
\setlength{\parindent}{0pt}
\setlength{\parskip}{0pt}

\titleformat{\section}
    {\sectiontitlefontsize\bfseries\color{sectiongray}}
    {}{0em}
    {}[\vspace{-14pt}{\color{applegray}\rule{\textwidth}{0.8pt}}]
\titlespacing{\section}{0pt}{8pt}{3pt}
\newcommand{\sectionbodyspace}{\vspace{-3pt}}
\newcommand{\sectiongap}{\vspace{-7pt}}
\newcommand{\entrygap}{\vspace{1pt}}
\newcommand{\projectgap}{\vspace{3pt}}

\newcommand{\resumeSubheading}[4]{
    \vspace{0pt}
    \begin{tabular*}{\textwidth}{l@{\extracolsep{\fill}}r}
        \textbf{\entrytitlefontsize\color{appledark}#1} \textit{\entrydetailsfontsize\color{textgray} · #3#4} & {\entrydetailsfontsize\color{applegray}#2} \\
    \end{tabular*}
    \vspace{-2pt}
}

\newcommand{\resumeItemListStart}{
    \begin{itemize}[
        leftmargin=1.5em,
        labelsep=0.5em,
        label={\color{applegray}\raisebox{0.12ex}{\scriptsize\faAngleRight}},
        itemsep=3pt,
        parsep=0pt,
        topsep=2pt
    ]
}
\newcommand{\resumeItemListEnd}{\end{itemize}\vspace{1pt}}
\newcommand{\skillItemListStart}{
    \begin{itemize}[
        leftmargin=1.5em,
        labelsep=0.5em,
        label={\color{applegray}\raisebox{0.12ex}{\scriptsize\faAngleRight}},
        itemsep=1pt,
        parsep=0pt,
        topsep=1pt
    ]
}

\newcommand{\headerLinkItem}[2]{
    \textcolor{textgray}{\linkfontsize#1}~\textcolor{appleblue}{\linkfontsize#2}
}

\newcommand{\resumeItem}[1]{\item\itemfontsize{#1}}

\begin{document}

\noindent
\begin{minipage}[t]{0.64\textwidth}
\vspace{0pt}
{\namefontsize\bfseries\color{appledark}刘龙昌}\par\vspace{4pt}
{\entrytitlefontsize\bfseries\color{textgray}软件开发工程师}\par\vspace{5pt}
{\contactfontsize\textcolor{textgray}{\faPhone}~\textcolor{textgray}{18357571400}\quad \textcolor{textgray}{\faEnvelope}~\textcolor{textgray}{liulch.cn@gmail.com}}\par
\end{minipage}\hfill
\begin{minipage}[t]{0.30\textwidth}
\vspace{0pt}
\raggedright
\vspace{6pt}
{\headerLinkItem{\faGithub}{\href{https://github.com/leochame}{github.com/leochame}}}\par\vspace{3pt}
{\headerLinkItem{\faLinkedin}{\href{https://www.linkedin.com/in/longchang-liu-770211387/}{LinkedIn / Longchang Liu}}}\par\vspace{3pt}
{\headerLinkItem{\faGlobe}{\href{https://leochame.github.io}{leochame.github.io}}}\par
\end{minipage}
\par\vspace{2pt}

\section{教育经历}
\sectionbodyspace
\begin{tabular*}{\textwidth}{l@{\extracolsep{\fill}}r}
    \textbf{\entrytitlefontsize\color{appledark}杭州电子科技大学} & {\entrydetailsfontsize\color{applegray}2022年08月 - 2026年06月} \\
\end{tabular*}
{\entrydetailsfontsize\color{textgray}计算机科学与技术 · 本科}\par
\sectiongap

\section{工作经历}
\sectionbodyspace
\resumeSubheading
{阿里巴巴}{2025年10月 - 2026年01月}{高德地图}{ · 实习 / Intern}
\entrygap
\resumeSubheading
{腾讯}{2025年04月 - 2025年09月}{广告营销服务线}{ · 实习 / Intern}
\sectiongap

\section{开源经历}
\sectionbodyspace
{\entrydetailsfontsize\color{textgray}\textbf{\color{appledark}LangChain4j}（12k+ Star）贡献者，Java 生态 AI 开发框架，参与社区问题修复与能力完善。}\par
\sectiongap

\section{项目经验}
\sectionbodyspace

\begin{tabular*}{\textwidth}{l@{\extracolsep{\fill}}r}
    \textbf{\smallentrytitlefontsize\color{appledark}阿里巴巴 - 高德地图 \entrydetailsfontsize\color{textgray} · } & {\entrydetailsfontsize\color{applegray}2025年10月 - 2026年01月} \\
\end{tabular*}
{\descfontsize\color{textgray} 负责“订单分析助手”与“AI 打车”的大模型工程化落地，提升复杂任务处理能力并完成统一 AI 能力接入。}\vspace{2pt}
\resumeItemListStart
\resumeItem{在订单分析助手项目中，设计并落地 Deep Research 功能，提升复杂任务拆解效果与长上下文场景下的处理表现。}
\resumeItem{参与 AI 打车服务冷启动，建设统一、稳定的 AI 调用接口，保障核心链路稳定接入。}
\resumeItemListEnd
\projectgap

\begin{tabular*}{\textwidth}{l@{\extracolsep{\fill}}r}
    \textbf{\smallentrytitlefontsize\color{appledark}腾讯 - 广告 ChatBI 系统 \entrydetailsfontsize\color{textgray} · } & {\entrydetailsfontsize\color{applegray}2025年04月 - 2025年09月} \\
\end{tabular*}
{\descfontsize\color{textgray} 参与腾讯广告 ChatBI 的 NL2Chart 核心链路建设，支撑自然语言请求到图表生成的转化。}\vspace{2pt}
\resumeItemListStart
    \resumeItem{为了保证源头语料库的质量，设计并实现语料自动化评审系统，推动核心链路召回正确率从 50\% 提升至 80\%。同时引入状态机，增强服务稳定性，并提供宕机恢复能力。}
\resumeItemListEnd
\sectiongap

\section{专业技能}
\vspace{0pt}
\skillItemListStart
    \resumeItem{\textbf{AI}：理解 Transformer 架构，了解 KV Cache、渐进式披露等基本原理与设计思想，熟悉 Agent 上下文管理以及常见的 Multi-Agent 协作范式，并辅以开源项目进行实践，具备实际落地经验}
    \resumeItem{\textbf{实践能力}：掌握 SQL 语句优化，MySQL 性能调优，线上故障排查、JVM 性能监控的基本方法}
    \resumeItem{\textbf{计算机基础}：具备扎实的操作系统与网络基础，深入理解内存管理（页表）、System Call、线程调度机制，以及 TCP/UDP 协议、拥塞控制等核心网络原理。}
    \resumeItem{\textbf{Java}：深入浅出 JVM 底层原理，熟悉类加载机制、内存模型以及主流垃圾收集器（G1/ZGC）}
    \resumeItem{\textbf{并发编程}：深入理解 Synchronized、Volatile、AQS 与线程池的实现原理}
\resumeItemListEnd

\end{document}
